\section{Introduction}\label{sec:introduction}

The Hilbert--Smith conjecture asks whether a locally compact group acting
faithfully and continuously on a connected topological manifold must be
a Lie group. An action is \emph{faithful} when its kernel is trivial:
no nonidentity element fixes every point. Faithfulness does not require
point stabilizers to be trivial. We prove the following statement.

\begin{theorem}[The Hilbert--Smith conjecture]\label{cor:lie}
Let \(n\geq1\). A locally compact second-countable Hausdorff group
acting faithfully and jointly continuously on a connected Hausdorff
second-countable topological \(n\)-manifold without boundary is a Lie
group.
\end{theorem}

This gives a positive resolution in every positive finite dimension.
The manifold may be noncompact, nonorientable, or nontriangulable;
fixed points and arbitrary stabilizers are allowed. In dimension zero,
a connected manifold is a point, so a faithful acting group is trivial
and the conclusion holds as well.

The mechanism is the exclusion of faithful actions of the additive
group of \(p\)-adic integers, for a prime \(p\),
\[
  \Zp=\varprojlim_k\Z/p^k\Z,
\]
with its profinite topology. The precise classical reduction from the
full group theorem to this case is recalled at the end of the proof.

\begin{theorem}[$p$-adic exclusion]\label{thm:main}
Let \(n\geq1\) be an integer and \(p\) a prime. Every jointly continuous
action of \(\Zp\) on a connected Hausdorff second-countable topological
\(n\)-manifold without boundary has nonzero kernel. Consequently the
action factors through a finite quotient of \(\Zp\).
\end{theorem}

For a manifold \(M\), equip \(\Homeo(M)\) with the compact-open
topology. Jointly continuous
actions are equivalently continuous homomorphisms
\(\Zp\to\Homeo(M)\). No dimension bound on the orbit space is assumed.

\Cref{cor:lie} also gives the dynamical consequence recorded in
\cite[Conjecture~1.4]{Pardon2013}. For a manifold $M$ as in
\Cref{cor:lie}, a homeomorphism $f$ is \emph{almost periodic} if the
cyclic subgroup $\{f^j:j\in\Z\}$ has compact closure $K$ in
$\Homeo(M)$. Some positive iterate of $f$ is then the time-one map
of a continuous real flow: there are an integer $m\geq1$ and a
continuous homomorphism $\Phi:\R\to\Homeo(M)$ with
$\Phi(1)=f^m$. Indeed, $K$ is a compact abelian group acting faithfully
and jointly continuously on $M$. Evaluation on a countable dense
subset of $M$ embeds $K$ into a countable product of copies of $M$,
so $K$ is second countable. By \Cref{cor:lie}, it is a compact abelian
Lie group. Its identity component $K^0$ is a torus of finite index,
so $f^m\in K^0$ for some $m\geq1$. The surjective exponential map
of this torus places $f^m$ on a one-parameter subgroup, giving $\Phi$.

\subsection{The problem and its predecessors}

Hilbert's fifth problem asked how far differentiability assumptions
could be removed from the theory of continuous transformation groups
\cite[Problem~5]{Hilbert1902}.
One central formulation concerns a group that is itself locally
Euclidean. Gleason and Montgomery--Zippin resolved that formulation in
1952 \cite{Gleason1952,MontgomeryZippin1952}, and Yamabe extended the
structure theory to general locally compact groups
\cite{Yamabe1953a,Yamabe1953b}.
Hilbert--Smith asks for the restrictions imposed by a faithful action
on a connected manifold. The passage to additive $p$-adic groups
combines this structure theory with rigidity of periodic
transformations. Lee gives a treatment of the reduction
\cite{Lee1997}; Pardon's introduction states the action-specific form
used here \cite[Introduction, p.~1]{Pardon2013}. The structure theory is
developed systematically in Montgomery--Zippin
\cite{MontgomeryZippin1955}.

The low-dimensional cases illustrate how much geometry may be needed to
control a merely continuous action.  The cases of one- and two-manifolds
were classical; Pardon gives a later account and proof of the
two-dimensional case, crediting Montgomery--Zippin
\cite{Pardon2019}.  Pardon proved the three-dimensional conjecture
\cite[Theorem~1.5]{Pardon2013}.  His local argument uses incompressible
surfaces in a suitable invariant open set and the resulting action on a
surface mapping class group.  In particular, it accommodates arbitrary
stabilizers without replacing the given action by a smooth action.

Another line of work controls the action by its regularity.
Bochner--Montgomery established the differentiable precedent
\cite{BochnerMontgomery1946}.
Repov\v{s}--\v{S}\v{c}epin proved the Lipschitz case on Riemannian
manifolds \cite[Theorem~1.1]{RepovsScepin1997}, and Maleshich obtained
an obstruction for H\"older actions on closed manifolds when a common
exponent exceeds $n/(n+2)$ \cite[Theorem~3.1]{Maleshich1997}.
Martin established the quasiconformal case on Riemannian manifolds
\cite[Theorems~1.2--1.3]{Martin1999}.
These results impose regularity on the transformations, rather than only
on the underlying manifold.

The metric arguments connect to the dimension-raising theory of
Yang and Bredon--Raymond--Williams
\cite{Yang1960,BredonRaymondWilliams1961}.
Their obstruction concerns integral cohomological dimension;
Raymond's subsequent analysis emphasizes the different behavior with
field coefficients \cite{Raymond1961}.
In the Lipschitz proof, an averaged invariant metric keeps Hausdorff
dimension small enough to contradict the required dimension increase
of the orbit space.
Mj extended this approach to uniformly quasisymmetric actions on
compact metric cohomology manifolds under Ahlfors regularity and
Hausdorff-dimension hypotheses, with applications to boundaries of
Poincar\'e duality hyperbolic groups
\cite[Theorem~2.24 and Corollary~2.25]{Mj2012}.

Shelukhin used Floer-theoretic methods to exclude nontrivial continuous
$p$-adic actions by homeomorphisms in the $C^0$-closure of the Hamiltonian
diffeomorphism group of a closed symplectic manifold on whose $\pi_2$
both the symplectic class and the first Chern class vanish
\cite[Theorem~A]{Shelukhin2024}.

Newman's rigidity theorem for periodic transformations
\cite{Newman1931} supplies the local input used here: a finite-order
homeomorphism of a connected manifold that fixes a nonempty open set
pointwise is the identity. For a modern statement and proof, see
\cite[author manuscript, pp.~2--3]{Pardon2019}.
Pardon's local reduction \cite[\S4.1]{Pardon2013} uses this rigidity to
find a faithful action of an open subgroup on an invariant open subset
of one coordinate chart. We use that reduction without imposing
regularity on the transformations or on the boundary of the chart subset.

\subsection{Duality and signatures}

The signature viewpoint has a separate line of antecedents.
Verdier developed duality for sheaves on locally compact spaces
\cite{Verdier1966}; we use its derived functor formulation in
Kashiwara--Schapira and Schapira
\cite{KashiwaraSchapira1990,SchapiraNotes}.
Balmer's triangular Witt groups package symmetric forms and their
neutral relations in categories with duality \cite{Balmer2000}.
For localization we use Balmer--Walter's published
Theorem~2.1, which removes the weak-cancellation hypothesis of the
earlier localization theorem \cite{BalmerWalter2002}.
Woolf applied this formalism both to perfect complexes of sheaves and
to constructible complexes. Using duality on perfect complexes, he
proved a cylinder homotopy relation for that theory and obtained a
generalized homology theory from constructible Witt groups on compact
polyhedra
\cite{Woolf2008}.
Ranicki--Weiss construct rational Pontryagin classes using a local
chain theory generated by continuous simplex maps, with cutoff by
subdivision, excision, and control of passage to direct summands
\cite[\S4 and Theorems~8.2--8.3]{RanickiWeiss2010}.
Their chain-level $L$-theory differs from the sheaf Witt groups used
here, but its generated-category strategy is a close antecedent.

Our categories are generated by proper images of finite polyhedra under
arbitrary continuous maps. Such images need not be constructible.
We verify that duality and the constructions in the cylinder relation
restrict to these specified generated categories. The additional
quotient comparison and uniform integral bound establish the fixed
signature lattice.
The comparison is integral before rationalization: a rational
isomorphism alone would not bound the possible denominators.

\subsection{The obstruction: equal parts in a fixed lattice}

The numerical contradiction is simple. Fix an odd integer
\(d\geq3\); for an action on an \(n\)-manifold, we will take \(d>n\).
We construct an abelian group \(E_d(S^d)\), an orientation
class \(u_d\), and a positive integer \(L_d\). For an integral class
\(x\), write \(\overline{x}=x\otimes1\) for its rational image. Then
\begin{equation}\label{intro:lattice}
 E_d(S^d)\otimes\Q=\Q\overline{u}_d,
 \qquad
 \im\bigl(E_d(S^d)\to E_d(S^d)\otimes\Q\bigr)
       \subseteq L_d^{-1}\Z\overline{u}_d.
\end{equation}
A faithful \(\Zp\)-action will produce, for every \(k\geq1\),
\(p^k\) integral classes in this group. Their rational images are equal
and sum to \(4\overline{u}_d\). Each image is therefore
\((4/p^k)\overline{u}_d\), which
contradicts \Cref{intro:lattice} as soon as \(p^k>4L_d\).
The order of these choices matters: \(L_d\) is fixed before \(k\).

Here ``integral'' refers to an element of the abelian Witt group before
it is tensored with \(\Q\); the sheaf coefficients are real throughout.
For a map from a compact polyhedron, its direct-image sheaf complex
records the cohomology of the fibers. A symmetric self-duality on that
complex supplies signature data, and the Witt group sets forms with
a lagrangian equal to zero. This makes the invariant additive
while retaining the integrality of ordinary signature.

Here is the group underlying this argument. For a finite polyhedron \(X\), let
\(\T(X)\) be the full thick subcategory of the bounded derived category
of real sheaves generated by
\[
 Rf_*\const{\R}{B},\qquad f:B\longrightarrow X,
 \qquad B\text{ a compact finite polyhedron},
\]
where \(f\) is continuous. Thus an object is obtained by finite sums,
shifts, cones, and direct summands from these proper direct images.
Their stalks need not be finite dimensional. Verdier duality nevertheless
restricts to this category, as we prove from simplex generators.
We put \(E_r(X)=W_r(\T(X))\), the Witt group of nonsingular symmetric
forms for the shifted duality \(\D_{X,r}=[-r]\D_X\), with neutral forms
set equal to zero. \Cref{sec:categories} gives the precise definitions,
the signs, and homotopy invariance.

To establish \Cref{intro:lattice}, we compare with the subcategory
generated by semialgebraic maps. On a suitably chosen oriented ball,
finite Witt relations can be evaluated by ordinary integral signature. Two
arguments transfer that integrality to the continuous category.
First, cutting off the finite polyhedral sources identifies quotient
morphisms with germs near a closed set. Second, a full dense
triangulated inclusion with duality, in which every object of the larger
category is a summand of one from the smaller category, admits a map
on Witt groups in
the reverse direction whose two composites are multiplication by
\(2\). The resulting exponent-two kernel and cokernel bound is a
consequence of Hornbostel--Schlichting's cofinality theorem
\cite[author preprint, Appendix~A, Theorem~A.2]{HornbostelSchlichting2004}.
We give the cutoff and explicit doubling constructions in
\Cref{sec:localization}; neither requires the images of arbitrary
continuous maps to be constructible.

Localization on a finite cover of the sphere then bounds both the
kernel and the cokernel of comparison by a fixed power of \(2\), in
every degree. This proves the lattice statement in
\Cref{thm:sphere-lattice}. Rational equivalence alone would lose the
denominator bound. The cutoff description and the integral bound for
dense inclusions are also useful independently of the group action;
they supply the additional comparison needed beyond the constructible
homology theory of \cite{Woolf2008}.

\subsection{How the action supplies the equal parts}

Suppose that a faithful action exists on an \(n\)-manifold. By Newman's
theorem, a small open subgroup \(G\cong\Zp\) acts faithfully and
preserves orientation on a connected invariant open subset
\(O\subset\R^n\) of one coordinate chart. Choose an odd integer
\(d>n\), with \(d\geq3\), and give
\[
 W=O\times\R^{d-n}\subset\R^d
\]
the product action, trivial in the added coordinates. We use the
one-point compactification \(W^+\), its orbit space \(Y=W^+/G\),
and the maps
\[
 S^d\xrightarrow{c}W^+\xrightarrow{q}Y\xrightarrow{t}X.
\]
Here \(c\) collapses \(S^d\setminus W\) to \(\infty\), and \(t\)
is a continuous map into a compact space \(X\) homeomorphic to a finite
polyhedron. We require the value
\(b=tq(\infty)\) to have a basis of contractible open neighborhoods
in \(X\). The composite \(tqc\) always has the same finite-polyhedron
source \(S^d\), as shown in \Cref{fig:local-test}.

\begin{figure}[!htbp]
\centering
\begin{tikzpicture}[
  >=Latex,
  every node/.style={font=\small},
  space/.style={draw,rounded corners=2pt,align=center,
    minimum height=12mm,inner sep=6pt},
  note/.style={font=\footnotesize,align=center,text=black!75}
]
\node[space] (source) at (0,0) {\(S^d\)};
\node[space] (compact) at (3.6,0) {\(W^+\)};
\node[space] (orbit) at (6.7,0) {\(Y=W^+/G\)};
\node[space] (test) at (10.0,0) {\(X\)};
\draw[->] (source) -- node[above]{\(c\)} (compact);
\draw[->] (compact) -- node[above]{\(q\)} (orbit);
\draw[->] (orbit) -- node[above]{\(t\)} (test);
\draw[->] (source.south) -- ++(0,-1.05)
  -| node[pos=.25,below]{\(f=tqc\)} (test.south);
\node[note,above=8mm of source] {compact polyhedral\\source};
\node[note,above=8mm of compact] {\(G\) acts on \(W^+\)\\and fixes \(\infty\)};
\node[note,above=8mm of orbit] {compact orbit\\space};
\node[note,above=8mm of test] {test space};
\node[note,below=17mm of compact] {
  \(c\) is the identity on \(W\) and collapses \(S^d\setminus W\) to \(\infty\).};
\end{tikzpicture}
\caption{A local action produces compact-source tests. Away from
\(b=tq(\infty)\), the source of \(f\) lies in the oriented manifold
\(W=O\times\R^{d-n}\). The restricted map
\(f^{-1}(X\setminus\{b\})\to X\setminus\{b\}\) is proper, with oriented
manifold source.}
\label{fig:local-test}
\end{figure}


Put \(f_t=tqc\) and \(A_t=Rf_{t*}\const{\C}{S^d}\).
Here \(\C\) is the positive real plane with pairing
\(B(z,w)=\operatorname{Re}(z\overline w)\). Pushing this coefficient
pairing together with the orientation of \(S^d\) gives the degree-\(d\)
form \(\alpha_t\) on \(A_t\).
Away from \(b\), the map \(f_t\) is proper from an oriented open
subset of \(W\). Identify the characters of \(G\) with the group
\(\Lambda=\bigcup_{j\geq0}\mu_{p^j}\) of roots of unity of $p$-power
order. A finite partition \(\Lambda=D_0\sqcup\cdots\sqcup D_{m-1}\)
gives orthogonal projectors \(e_i\) on
\(A_t|_{X\setminus\{b\}}\) onto the corresponding character parts.
A projector need not extend to a summand over \(b\). Choose a
contractible open neighborhood \(U\) of \(b\). The projector's germ near
\(X\setminus U\) defines an endomorphism of \(A_t\) in the quotient
by objects compactly supported in \(U\).

For a self-adjoint projector \(e_i\), the operator \(2e_i-1\) is a
self-adjoint involution. We therefore have the following sum of forms
on the existing object:
\[
 [A_t,\alpha_t]+[A_t,\alpha_t(2e_i-1)].
\]
An integral localization
isomorphism lifts it to a class \(z_i(t)\in E_d(X)\). An explicit
matrix isometry gives
\(\sum_i z_i(t)=2[A_t,\alpha_t]\), without taking the images of
quotient idempotents.

Averaging a map supported in a ball gives a test \(a:Y\to S^d\)
with \(\deg(aqc)=1\). For each \(k\), faithfulness gives a nonempty
open region \(V\subset W/G\) over which
\(W/(p^kG)\to W/G\) consists of \(p^k\) labeled sheets.
Serre's results on odd spheres provide a map \(D_s:S^d\to S^d\)
of degree \(s>0\) such that \(D_s a\) is homotopic to a map constant
outside \(V\) \cite{Serre1951,Serre1953}. There a locally constant coefficient
isometry permutes \(p^k\) chosen character parts cyclically.
Covariance and homotopy invariance equate the rational images
\(s\overline{z}_i(a)\). Cancelling the nonzero rational scalar \(s\)
equates the \(\overline{z}_i(a)\), while the original integral classes
stay in the same fixed lattice. Their rational images sum to
\(4\overline{u}_d\): one factor of \(2\) comes from the real coefficient
plane and the other from the involution construction. The fixed lattice
gives the contradiction.

Only the actual manifold sources and the test spaces carry Verdier
duality in this argument. The compact orbit spaces enter through
ordinary sheaf direct images, averaging, and \v{C}ech cohomology.
No dimension bound on the orbit spaces or regularity assumption on
the boundary of \(O\) is used.
The odd stabilization is the only dimensional choice. In particular,
the proof treats both odd congruence classes of \(d\) modulo \(4\).

\Cref{sec:categories,sec:localization,sec:lattice} construct and bound
the signature invariant. \Cref{sec:local-tests} prepares the local
tests, including the odd-sphere argument. \Cref{sec:characters}
constructs the integral character classes and proves
\Cref{thm:main,cor:lie}.
